% META: {"id":"1NSI-LANGAGE-RE-C4-CORRIGE","chapitre":"1NSI-LANGAGE","type_objet":"corrige","exercice_ref":"1NSI-LANGAGE-RE-C4","capacites":["P-LANG-04"],"status":"approved","origin":"needs_review","status_history":["needs_review","audited","accepted","approved"],"release_acceptance":"RELEASE_OWNER_BATCH_ACCEPTANCE","acceptance_closure_digest":"sha256:9b3ccf9a81c5520fbb7e03b7b2d3e7bf2057a02834b6d908bf3be5f49f3b553f","acceptance_receipt_digest":"sha256:4fad86f0282e0fa4e0419cca1ad6379ae7af5a280c04a4a90880f90a1c044242"}

\begin{corrige}{1NSI-LANG-RE-C4-EX1}
\textbf{1.} \lstinline{minimum_bugue([5, 3, 8])} renvoie \lstinline{0}, alors que le
minimum réel de la liste est $3$.

\textbf{2.} L'accumulateur \lstinline{mini} est initialisé à $0$, qui est \textbf{inférieur
à tous les éléments} de la liste \lstinline{[5, 3, 8]} (tous positifs et $\geqslant3$).
Aucun élément ne vérifie donc \lstinline{x < mini}, et \lstinline{mini} reste égal à $0$
jusqu'à la fin : ce $0$, qui n'appartient même pas à la liste, est renvoyé à tort.

\textbf{3.}
\textbf{Précondition : la liste est non vide.} Un test explicite refuse donc
\lstinline{[]} en déclenchant une \lstinline{ValueError}.

% PYTHON-SOURCE: code/minimum.py
\begin{python}
def minimum(liste):
    """Renvoie le plus petit element.

    Precondition : liste est non vide.
    """
    if len(liste) == 0:
        raise ValueError("liste doit etre non vide")
    mini = liste[0]
    for x in liste[1:]:
        if x < mini:
            mini = x
    return mini
\end{python}
En initialisant \lstinline{mini} avec le \textbf{premier élément} de la liste (et non une
constante arbitraire), la fonction renvoie bien $3$ pour \lstinline{[5, 3, 8]}.
\end{corrige}

% BEGIN-VERIFY
% def minimum(liste):
%     """Renvoie le plus petit element.
%
%     Precondition : liste est non vide.
%     """
%     if len(liste) == 0:
%         raise ValueError("liste doit etre non vide")
%     mini = liste[0]
%     for x in liste[1:]:
%         if x < mini:
%             mini = x
%     return mini
% assert minimum([5, 3, 8]) == 3
% try:
%     minimum([])
% except ValueError as erreur:
%     assert str(erreur) == "liste doit etre non vide"
% else:
%     raise AssertionError("minimum doit refuser une liste vide")
% END-VERIFY
